% Part: first-order-logic
% Chapter: syntax-and-semantics
% Section: discussion

\documentclass[../../../include/open-logic-section]{subfiles}

\begin{document}

\olfileid{fol}{dis}{exi}

\olsection{Existence in first-order logic}

\paragraph{To be is to be something} In first-order logic we can express the idea that something exist. Namely, to be it to be something. For instance, we express that $a$ exists thus:

$$\lexists x(\eq[a][x])$$$

\paragraph{Some consequences of first-order logic}

Can you see that the following holds in first-order logic? 

\begin{enumerate}
	\item $\Entails \lforall (\eq[x][x])$.
	\item For any constant $c$, we have: $\Entails \lexists (\eq[c][x])$.
	\item $\Entails \lexists (\eq[x][x])$.
\end{enumerate} 

Let's look ahead to the combination of first-order logic and modal logic (quantified modal logic). In "normal" modal logic, we have the principle that everything that is logically valid is necessary:

\begin{defn}[Necessitation]
If $\Entails !A$ then $\Entails \Box !A$ 
\end{defn}

Combining our result above with Necessitation, we obtain: 

\begin{enumerate}
	\item $\Entails \Box \lforall (\eq[x][x])$.
	\item For any constant $c$, we have: $\Entails \Box \lexists (\eq[c][x])$.
	\item $\Entails \Box \lexists (\eq[x][x])$.
\end{enumerate} 

Which of these do you find acceptable? Inacceptable?

\end{document}
